\chapter{Appendix}
\label{ch:appendix}
\raggedbottom

\section{Regression Diagnostics}

This appendix reports only supporting material that is useful for checking the main analysis without repeating the figures already discussed in the thesis chapters. Table~\ref{tab:appendix_full_pooled_regression} gives the full pooled baseline regression, including the year indicators omitted from the compact main table. The omitted year is 2015.

\begin{table}[H]
\centering
\scriptsize
\caption{Full pooled baseline regression with year indicators}
\label{tab:appendix_full_pooled_regression}
\begin{tabular}{@{}lrrrr@{}}
\toprule
Variable & Estimate & Std. error & \(t\) & \(p\) \\
\midrule
Intercept & 7.390 & 6.813 & 1.085 & 0.278 \\
Year 2016 & -0.033 & 0.241 & -0.135 & 0.892 \\
Year 2017 & -0.428 & 0.416 & -1.028 & 0.304 \\
Year 2018 & 0.000 & 0.492 & 0.000 & 1.000 \\
Year 2019 & 0.025 & 0.559 & 0.046 & 0.964 \\
Year 2020 & -0.136 & 0.470 & -0.289 & 0.772 \\
Year 2021 & -0.171 & 0.553 & -0.308 & 0.758 \\
Year 2022 & 0.242 & 0.733 & 0.330 & 0.742 \\
Year 2023 & 0.804 & 0.772 & 1.042 & 0.297 \\
Year 2024 & 0.803 & 0.630 & 1.274 & 0.203 \\
Log GDP per capita & -0.546 & 0.593 & -0.920 & 0.358 \\
Unemployment rate & 0.160 & 0.186 & 0.861 & 0.389 \\
Poverty or social exclusion & 0.062 & 0.065 & 0.955 & 0.340 \\
Government health expenditure & 0.159 & 0.190 & 0.835 & 0.404 \\
Compulsory health financing & -0.401 & 0.243 & -1.652 & 0.099 \\
\bottomrule
\end{tabular}
\par\vspace{0.3em}
\begin{minipage}{0.96\textwidth}
\footnotesize Notes: This is the full pooled baseline specification corresponding to Table~\ref{tab:regression_main_results}, including all year-indicator coefficients. The omitted year is 2015. Standard errors are country-clustered; exact p-values are shown instead of significance stars.
\end{minipage}

\end{table}

Table~\ref{tab:appendix_regression_vif_diagnostics} reports variance inflation factors for the baseline covariates. The diagnostics are included to document the degree of collinearity among the main regressors used in the complete-case model.

\begin{table}[H]
\centering
\small
\caption{Variance inflation factors for baseline covariates}
\label{tab:appendix_regression_vif_diagnostics}
\begin{tabular}{@{}lr@{}}
\toprule
Covariate & VIF \\
\midrule
Log GDP per capita & 1.92 \\
Unemployment rate & 1.28 \\
Poverty or social exclusion & 2.00 \\
Government health expenditure & 2.31 \\
Compulsory health financing & 2.82 \\
\bottomrule
\end{tabular}
\par\vspace{0.3em}
\begin{minipage}{0.86\textwidth}
\footnotesize Notes: VIF is the variance inflation factor calculated for the five baseline covariates in the complete-case pooled regression sample. Values above 10 are commonly treated as indicating serious multicollinearity.
\end{minipage}

\end{table}

\section{Missingness and Prediction Sensitivities}

Table~\ref{tab:appendix_mi_variant_summary} documents the PMM multiple-imputation variants used to assess whether the main attenuation pattern depends on including the outcome, country-structure representation, lagged predictors, or the number of imputations.

\begin{table}[H]
\centering
\small
\caption{Multiple-imputation variants for the poverty coefficient}
\label{tab:appendix_mi_variant_summary}
\begin{tabular}{@{}lrrrr@{}}
\toprule
Variant & Coef. & SE & p-value & \(M\) \\
\midrule
Baseline PMM MI & 0.0139 & 0.0434 & 0.748 & 30 \\
No-outcome PMM MI & -0.0308 & 0.0492 & 0.532 & 30 \\
Country-FE PMM MI & 0.2516 & 0.0896 & 0.005 & 30 \\
Lagged-predictor PMM MI & -0.0472 & 0.0566 & 0.405 & 30 \\
50-imputation PMM check & 0.0122 & 0.0432 & 0.778 & 50 \\
\bottomrule
\end{tabular}

\end{table}

Table~\ref{tab:appendix_mnar_sensitivity_results} documents MNAR delta-shift scenarios. These scenarios shift only originally missing values after PMM imputation and are interpreted as stress tests rather than corrected estimates.

\begin{table}[H]
\centering
\scriptsize
\caption{MNAR delta-shift sensitivity for the poverty coefficient}
\label{tab:appendix_mnar_sensitivity_results}
\resizebox{\textwidth}{!}{%
\begin{tabular}{@{}lrrrl@{}}
\toprule
Scenario & Coef. [95\% CI] & p-value & \(M\) & Interpretation \\
\midrule
GDP -10\% & 0.011 [-0.085, 0.107] & 0.823 & 30 & MNAR stress test \\
GDP -20\% & 0.012 [-0.085, 0.108] & 0.811 & 30 & MNAR stress test \\
GDP -5\% & 0.011 [-0.086, 0.107] & 0.828 & 30 & MNAR stress test \\
Optimistic combo & 0.010 [-0.086, 0.106] & 0.837 & 30 & MNAR stress test \\
Pessimistic combo & 0.010 [-0.086, 0.107] & 0.831 & 30 & MNAR stress test \\
Poverty +10 pp & 0.007 [-0.089, 0.103] & 0.890 & 30 & MNAR stress test \\
Poverty +2 pp & 0.010 [-0.086, 0.107] & 0.834 & 30 & MNAR stress test \\
Poverty +5 pp & 0.010 [-0.086, 0.106] & 0.845 & 30 & MNAR stress test \\
Unemployment +2 pp & 0.011 [-0.085, 0.107] & 0.827 & 30 & MNAR stress test \\
Unemployment +5 pp & 0.011 [-0.084, 0.107] & 0.817 & 30 & MNAR stress test \\
\bottomrule
\end{tabular}%
}

\vspace{0.2em}
\begin{minipage}{0.96\textwidth}\scriptsize Notes: MNAR scenarios apply documented shifts only to originally missing values after the PMM MAR imputation. They are stress tests for model dependence and should not be read as corrected estimates.\end{minipage}

\end{table}

Table~\ref{tab:appendix_ml_naive_baselines} reports the naive prediction baselines used to decide whether flexible models add useful one-year-ahead predictive value.

\begin{table}[H]
\centering
\small
\caption{Naive prediction baselines}
\label{tab:appendix_ml_naive_baselines}
\begin{tabular}{@{}rrrr@{}}
\toprule
Model & MAE & RMSE & R2 \\
\midrule
Previous-year outcome & 0.587 & 1.006 & 0.850 \\
AR(1) linear & 0.661 & 1.101 & 0.820 \\
Ridge with lagged predictors & 0.695 & 1.115 & 0.816 \\
Last observed country value & 0.718 & 1.187 & 0.791 \\
Rolling country mean & 0.796 & 1.357 & 0.727 \\
Country mean & 0.855 & 1.340 & 0.734 \\
Last training-year mean & 1.815 & 2.613 & -0.013 \\
Global mean & 1.885 & 2.596 & -0.000 \\
\bottomrule
\end{tabular}

\end{table}

\section{Data Sources and Reproducibility}

Table~\ref{tab:appendix_eurostat_sources} lists the Eurostat tables and filters used to construct the datasets. The primary outcome is hlth\_silc\_08, while hlth\_silc\_08b is used only as the robustness outcome because its denominator differs from the main population-level indicator.

\begin{table}[H]
\centering
\scriptsize
\caption{Eurostat source tables, filters, and extraction dates}
\label{tab:appendix_eurostat_sources}
\resizebox{\textwidth}{!}{%
\begin{tabular}{@{}>{\raggedright\arraybackslash}p{0.15\textwidth}>{\raggedright\arraybackslash}p{0.23\textwidth}>{\raggedright\arraybackslash}p{0.45\textwidth}r@{}}
\toprule
Code & Use & Filters & Accessed \\
\midrule
\texttt{hlth\_silc\_08} & Medical population denominator & unit=PC; quantile=TOTAL; reason=TXP\_TFAR\_WLIST; age=Y\_GE16; sex=T; years 2008--2025 & 2026-06-23 \\
\texttt{hlth\_silc\_08b} & Medical need denominator & unit=PC; rskpovth=TOTAL; reason=TXP\_TFAR\_WLIST; age=Y\_GE16; sex=T; years 2021--2025 & 2026-06-23 \\
\texttt{hlth\_silc\_09} & Dental population denominator & unit=PC; quantile=TOTAL; reason=TXP\_TFAR\_WLIST; age=Y\_GE16; sex=T; years 2008--2025 & 2026-06-23 \\
\texttt{hlth\_silc\_09b} & Dental need denominator & unit=PC; rskpovth=TOTAL; reason=TXP\_TFAR\_WLIST; age=Y\_GE16; sex=T; years 2021--2025 & 2026-06-23 \\
\texttt{demo\_pjan} & Population weights & freq=A; unit=NR; age=TOTAL; sex=T; years 2008--2025 & 2026-05-21 \\
\texttt{nama\_10\_pc} & GDP per capita & unit=CP\_EUR\_HAB; na\_item=B1GQ; years 2008--2025 & 2026-05-14 \\
\texttt{une\_rt\_a} & Unemployment rate & unit=PC\_ACT; sex=T; age=Y15-74; years 2008--2025 & 2026-05-14 \\
\texttt{ilc\_peps01n} & Poverty or social exclusion & unit=PC; sex=T; age=TOTAL; years 2008--2025 & 2026-05-14 \\
\texttt{gov\_10a\_exp} & Government health expenditure & unit=PC\_GDP; sector=S13; cofog99=GF07; na\_item=TE; years 2008--2025 & 2026-05-14 \\
\texttt{hlth\_sha11\_hf} & Compulsory health financing & unit=PC\_GDP; icha11\_hf=HF1; years 2008--2025 & 2026-05-14 \\
\texttt{hlth\_sha11\_hf} & Out-of-pocket spending & unit=PC\_CHE; icha11\_hf=HF3; years 2008--2025 & 2026-05-14 \\
\texttt{hlth\_rs\_prs2} & Physicians per 100,000 & unit=P\_HTHAB; wstatus=PRACT; med\_spec=PHYS; years 2008--2025 & 2026-05-14 \\
\texttt{hlth\_rs\_bds1} & Hospital beds per 100,000 & unit=P\_HTHAB; facility=HBEDT; hlthcare=TOTAL; years 2008--2025 & 2026-05-14 \\
\texttt{ilc\_di12} & Gini coefficient & age=TOTAL; statinfo=GINI\_HND; years 2008--2025 & 2026-05-14 \\
\texttt{une\_ltu\_a} & Long-term unemployment & indic\_em=LTU; age=Y15-74; sex=T; unit=PC\_ACT; years 2008--2025 & 2026-05-14 \\
\bottomrule
\end{tabular}%
}

\end{table}

Table~\ref{tab:appendix_software_versions} records the main software environment used for data processing, modelling, and figure generation.

\begin{table}[H]
\centering
\small
\caption{Python and library versions}
\label{tab:appendix_software_versions}
\begin{tabular}{@{}ll@{}}
\toprule
Software & Version \\
\midrule
Python & 3.13.12 \\
pandas & 3.0.3 \\
scikit-learn & 1.7.1 \\
statsmodels & 0.14.6 \\
matplotlib & 3.10.8 \\
seaborn & 0.13.2 \\
\bottomrule
\end{tabular}

\end{table}

Table~\ref{tab:processed_data_dictionary_summary} summarises the processed data files submitted with the thesis package. The full column-level data dictionary is written to \texttt{outputs/processed\_data\_dictionary.csv}.

\begingroup
\scriptsize
\begin{longtable}{@{}>{\raggedright\arraybackslash}p{0.31\textwidth}rr>{\raggedright\arraybackslash}p{0.33\textwidth}@{}}
\caption{Processed data dictionary summary}\label{tab:processed_data_dictionary_summary}\\
\toprule
File & Rows & Columns & Key columns \\
\midrule
\endfirsthead
\toprule
File & Rows & Columns & Key columns \\
\midrule
\endhead
citizenship.csv & 1612 & 4 & ctry, year, unmet \\
outcome.csv & 608 & 4 & geo, year, unmet \\
outcome\_08b.csv & 154 & 4 & geo, year \\
gini.csv & 410 & 3 & geo, year \\
beds.csv & 570 & 3 & geo, year \\
lt\_unemp.csv & 572 & 3 & geo, year \\
oop\_share.csv & 475 & 3 & geo, year \\
physicians.csv & 459 & 3 & geo, year \\
population.csv & 796 & 4 & geo, year, pop \\
ml\_5a.csv & 230 & 22 & geo, year, unmet \\
ml\_5a\_splits.csv & 230 & 23 & geo, year, unmet \\
ml\_time.csv & 236 & 15 & geo, year, unmet \\
multi\_outcome\_unmet... & 2740 & 11 & geo, year \\
panel.csv & 608 & 25 & geo, year, unmet, pop \\
imp\_1.csv & 608 & 25 & geo, year, unmet, pop \\
imp\_10.csv & 608 & 25 & geo, year, unmet, pop \\
imp\_2.csv & 608 & 25 & geo, year, unmet, pop \\
imp\_3.csv & 608 & 25 & geo, year, unmet, pop \\
imp\_4.csv & 608 & 25 & geo, year, unmet, pop \\
imp\_5.csv & 608 & 25 & geo, year, unmet, pop \\
imp\_6.csv & 608 & 25 & geo, year, unmet, pop \\
imp\_7.csv & 608 & 25 & geo, year, unmet, pop \\
imp\_8.csv & 608 & 25 & geo, year, unmet, pop \\
imp\_9.csv & 608 & 25 & geo, year, unmet, pop \\
skeleton.csv & 608 & 9 & geo, year, unmet \\
controls.csv & 667 & 7 & geo, year \\
\bottomrule
\end{longtable}
\endgroup


Table~\ref{tab:output_generation_map} gives a compact map from generated thesis tables and figures to their scripts and main inputs. Full paths and exact source filenames are written to \texttt{outputs/table\_figure\_generation\_map.csv}.

\begingroup
\tiny
\begin{longtable}{@{}p{0.11\textwidth}p{0.39\textwidth}p{0.13\textwidth}p{0.11\textwidth}@{}}
\caption{Table and figure generation map}\label{tab:output_generation_map}\\
\toprule
Use & Output & Script & Input \\
\midrule
\endfirsthead
\toprule
Use & Output & Script & Input \\
\midrule
\endhead
App. & 08b\_minus\_08\_differen... & extra & C/08b \\
App. & 08b\_vs\_08\_scatter.pdf & extra & C/08b \\
App. & citizenship\_gap\_corre... & extra & C/08b \\
App. & citizenship\_unmet\_nee... & extra & C/08b \\
App. & citizenship\_unmet\_nee... & extra & C/08b \\
App. & citizenship\_unmet\_nee... & extra & C/08b \\
Ch. 4 & complete\_case\_attriti... & full & Panel \\
Ch. 2 & conceptual\_selection\_... & diagram & Concept \\
App. & convergence\_sigma.pdf & extra & C/08b \\
Ch. 4 & descriptive\_country\_r... & full & Panel \\
App. & descriptive\_country\_r... & full & Panel \\
Ch. 4 & descriptive\_country\_y... & full & Panel \\
Ch. 4 & descriptive\_estimand\_... & full & Panel \\
App. & descriptive\_eu\_trend.pdf & full & Panel \\
App. & descriptive\_trend\_eur... & full & Panel \\
App. & descriptive\_trend\_wit... & extra & C/08b \\
App. & leave\_one\_country\_out... & missing & Panel \\
App. & leave\_one\_year\_out\_po... & missing & Panel \\
Ch. 5 & mi\_coefficient\_distri... & missing & Panel \\
App. & mi\_imputed\_share\_heat... & missing & Panel \\
App. & mi\_imputed\_share\_heat... & missing & Panel \\
App. & mi\_imputed\_share\_heat... & missing & Panel \\
App. & mi\_imputed\_share\_heat... & missing & Panel \\
App. & mi\_imputed\_share\_heat... & missing & Panel \\
App. & mi\_imputed\_share\_heat... & missing & Panel \\
App. & mi\_imputed\_share\_heat... & missing & Panel \\
App. & mi\_imputed\_share\_heat... & missing & Panel \\
App. & mi\_imputed\_share\_heat... & missing & Panel \\
Ch. 5 & mi\_observed\_vs\_impute... & missing & Panel \\
App. & mi\_observed\_vs\_impute... & missing & Panel \\
App. & mi\_observed\_vs\_impute... & missing & Panel \\
App. & mi\_observed\_vs\_impute... & missing & Panel \\
App. & mi\_observed\_vs\_impute... & missing & Panel \\
App. & mi\_observed\_vs\_impute... & missing & Panel \\
App. & mi\_observed\_vs\_impute... & missing & Panel \\
App. & mi\_observed\_vs\_impute... & missing & Panel \\
App. & mi\_observed\_vs\_impute... & missing & Panel \\
Ch. 3; Ch. 4 & missingness\_heatmap.pdf & full & Panel \\
App. & missingness\_heatmap\_c... & full & Panel \\
Ch. 3; Ch. 4 & missingness\_heatmap\_f... & full & Panel \\
App. & missingness\_heatmap\_o... & full & Panel \\
App. & missingness\_heatmap\_p... & full & Panel \\
Ch. 3; Ch. 4 & missingness\_heatmap\_s... & full & Panel \\
Ch. 5 & missingness\_inclusion... & full & Panel \\
App. & ml\_actual\_vs\_predicte... & ml & ML \\
App. & ml\_actual\_vs\_predicte... & ml & ML \\
App. & ml\_actual\_vs\_predicte... & ml & ML \\
App. & ml\_calibration\_plot.pdf & ml & ML \\
App. & ml\_expanding\_window\_v... & ml & ML \\
App. & ml\_importance\_compari... & ml & ML \\
App. & ml\_importance\_stabili... & ml & ML \\
App. & ml\_learning\_curves.pdf & ml & ML \\
App. & ml\_permutation\_vs\_imp... & ml & ML \\
App. & ml\_residuals\_by\_count... & ml & ML \\
App. & ml\_residuals\_by\_count... & ml & ML \\
App. & ml\_residuals\_by\_count... & ml & ML \\
App. & ml\_residuals\_vs\_predi... & ml & ML \\
App. & ml\_residuals\_vs\_predi... & ml & ML \\
App. & ml\_residuals\_vs\_predi... & ml & ML \\
App. & ml\_tree\_feature\_impor... & ml & ML \\
App. & ml\_validation\_tuning\_... & ml & ML \\
Ch. 5 & mnar\_tipping\_point\_cu... & missing & Panel \\
App. & model\_correlations\_he... & full & Panel \\
Ch. 5 & model\_ladder\_coeffici... & full & Panel \\
Ch. 4 & multi\_outcome\_monitor... & full & Panel \\
Ch. 5 & multi\_outcome\_primary... & full & Panel \\
Ch. 5 & multi\_outcome\_retenti... & full & Panel \\
Ch. 5 & outcome\_estimand\_coef... & full & Panel \\
App. & outcome\_missingness\_h... & full & Panel \\
App. & panel\_correlations\_he... & full & Panel \\
App. & panel\_heterogeneity.pdf & full & Panel \\
App. & panel\_variable\_distri... & full & Panel \\
Ch. 3 & population\_weighted\_t... & full & Panel \\
App. & poverty\_robustness\_fo... & missing & Panel \\
App. & regression\_fe\_coeffic... & missing & Panel \\
App. & regression\_robustness... & full & Panel \\
App. & scatter\_unmet\_vs\_comp... & full & Panel \\
App. & scatter\_unmet\_vs\_gdp\_... & full & Panel \\
App. & scatter\_unmet\_vs\_gini... & full & Panel \\
App. & scatter\_unmet\_vs\_gove... & full & Panel \\
App. & scatter\_unmet\_vs\_hosp... & full & Panel \\
App. & scatter\_unmet\_vs\_log\_... & full & Panel \\
App. & scatter\_unmet\_vs\_long... & full & Panel \\
App. & scatter\_unmet\_vs\_oop\_... & full & Panel \\
App. & scatter\_unmet\_vs\_phys... & full & Panel \\
App. & scatter\_unmet\_vs\_pove... & missing & Panel \\
App. & scatter\_unmet\_vs\_unem... & full & Panel \\
Ch. 5 & simulation\_bias\_by\_me... & full & Panel \\
Ch. 5 & simulation\_target\_dis... & full & Panel \\
App. & tree\_feature\_importan... & full & Panel \\
App. & unmet\_need\_trends\_all... & full & Panel \\
App. & unmet\_need\_vs\_gdp\_per... & full & Panel \\
App. & unmet\_need\_vs\_poverty... & missing & Panel \\
App. & unmet\_need\_vs\_unemplo... & full & Panel \\
App. & appendix\_full\_pooled\_... & full & Panel \\
Ch. 4 & attrition\_by\_country\_... & full & Panel \\
Ch. 4 & attrition\_by\_country\_... & full & Panel \\
Ch. 4 & attrition\_by\_covariat... & full & Panel \\
App. & attrition\_by\_year.tex & full & Panel \\
App. & citizenship\_gap\_clust... & extra & C/08b \\
App. & citizenship\_gap\_regre... & extra & C/08b \\
App. & citizenship\_gap\_table... & extra & C/08b \\
Ch. 4 & complete\_case\_attriti... & full & Panel \\
Ch. 5 & complete\_case\_selecti... & missing & Panel \\
App. & country\_coverage.tex & full & Panel \\
Ch. 3 & country\_universe.tex & missing & Panel \\
App. & descriptive\_statistic... & full & Panel \\
Ch. 5 & estimand\_framework.tex & missing & Panel \\
App. & eurostat\_sources\_appe... & summary & Gen. \\
Ch. 7 & failure\_modes\_public\_... & full & Panel \\
Ch. 5 & fe\_fit\_diagnostics.tex & missing & Panel \\
Ch. 5 & fe\_inference\_comparis... & missing & Panel \\
Ch. 7 & final\_classification\_... & full & Panel \\
Ch. 4 & included\_vs\_excluded\_... & full & Panel \\
Ch. 4 & leave\_one\_variable\_ou... & missing & Panel \\
Ch. 5 & main\_sensitivity\_ladd... & missing & Panel \\
Ch. 5 & mi\_diagnostics\_summar... & missing & Panel \\
Ch. 5 & mi\_pmm\_donor\_diagnost... & missing & Panel \\
App.; Ch. 5 & mi\_variant\_summary.tex & missing & Panel \\
Ch. 3; Ch. 4 & missingness\_compariso... & full & Panel \\
Ch. 5 & missingness\_robustnes... & missing & Panel \\
Ch. 6 & ml\_country\_holdout\_su... & ml & ML \\
App. & ml\_hyperparameter\_gri... & ml & ML \\
Ch. 6 & ml\_model\_performance\_... & ml & ML \\
App.; Ch. 6 & ml\_naive\_baselines.tex & ml & ML \\
App. & ml\_paired\_tree\_mae\_di... & ml & ML \\
App. & ml\_test\_mae\_bootstrap... & ml & ML \\
App. & ml\_time\_aware\_fold\_re... & ml & ML \\
App.; Ch. 5 & mnar\_sensitivity\_resu... & missing & Panel \\
Ch. 5 & mnar\_tipping\_point\_re... & missing & Panel \\
Ch. 5 & model\_ladder\_results.tex & full & Panel \\
Ch. 5 & multi\_outcome\_attriti... & full & Panel \\
App. & multi\_outcome\_country... & full & Panel \\
Ch. 4 & multi\_outcome\_coverag... & full & Panel \\
Ch. 5 & multi\_outcome\_include... & full & Panel \\
Ch. 4 & multi\_outcome\_monitor... & full & Panel \\
Ch. 5 & multi\_outcome\_primary... & full & Panel \\
App. & multiple\_imputation\_s... & full & Panel \\
Ch. 5 & outcome\_estimand\_coef... & full & Panel \\
Ch. 5 & outcome\_failure\_class... & full & Panel \\
App. & output\_generation\_map... & summary & Gen. \\
App. & poverty\_robustness\_su... & missing & Panel \\
App. & processed\_data\_dictio... & summary & Gen. \\
Ch. 3; Ch. 4 & public\_data\_estimand\_... & full & Panel \\
App. & reduced\_covariate\_sen... & full & Panel \\
App. & regression\_country\_gr... & full & Panel \\
App. & regression\_crisis\_exc... & full & Panel \\
App. & regression\_lag\_robust... & full & Panel \\
Ch. 5 & regression\_main\_resul... & full & Panel \\
App. & regression\_multilevel... & full & Panel \\
App. & regression\_robustness... & full & Panel \\
App. & regression\_sensitivit... & full & Panel \\
App. & regression\_sensitivit... & full & Panel \\
App. & regression\_vif.tex & full & Panel \\
App. & regression\_vif\_diagno... & full & Panel \\
Ch. 3 & robustness\_08b\_result... & extra & C/08b \\
App. & simulation\_complete\_c... & full & Panel \\
Ch. 5 & simulation\_main\_summa... & full & Panel \\
App. & software\_versions.tex & summary & Gen. \\
Ch. 7 & target\_population\_dri... & missing & Panel \\
Ch. 5 & target\_population\_sen... & missing & Panel \\
Ch. 5 & universe\_sensitivity\_... & full & Panel \\
Ch. 3 & variable\_summary.tex & full & Panel \\
\bottomrule
\end{longtable}
\endgroup


The computational workflow is fully scripted in the local project directory. Raw Eurostat JSON downloads are stored in \texttt{data/raw}; processed country-year files are stored in \texttt{data/processed}; generated CSV diagnostics are stored in \texttt{outputs}; generated LaTeX tables are stored in \texttt{tables}; and generated figures are stored in \texttt{figures}. The command order is:
\begin{verbatim}
rtk python src/run_full_thesis_analysis.py
rtk python src/generate_conceptual_selection_diagram.py
rtk python src/run_missingness_robustness.py
rtk python src/run_step5_ml_analysis.py
rtk python src/run_step6_additional_analysis.py
rtk python src/run_step7_master_summary.py
\end{verbatim}
The raw-data download workflow is part of the first command: it calls the Eurostat API for the main outcome, robustness outcome, citizenship table, population weights, and control/feature covariates using the filters listed in Table~\ref{tab:appendix_eurostat_sources}. The main table and figure generation map is:
\begin{itemize}
    \item Step 1 builds the panel, descriptive figures, baseline regressions, population-weighted trends, and source register;
    \item the diagram script creates the conceptual selection and estimand figure;
    \item the missingness script creates target-population sensitivity, IPW, MI, and poverty-robustness outputs;
    \item the prediction script creates one-year-ahead baselines, reduced benchmark-model comparisons, and country-holdout validation;
    \item the additional-analysis script creates citizenship paired-gap checks and \texttt{hlth\_silc\_08b} denominator robustness;
    \item the summary script creates the manifest, software versions, command order, hashes, and output checks.
\end{itemize}
All stochastic analyses use random seed 42 unless explicitly stated otherwise. The reproducibility manifest records Python and package versions, input and script hashes, output hashes, command order, the population-weighting source, and the machine-learning hyperparameter grids. The submitted package also includes raw-download instructions, an output hash manifest, and a reproducibility archive under \texttt{outputs}. The thesis does not cite a public code repository; reproducibility is documented through the submitted local scripts, raw downloads, processed files, generated outputs, source register, and manifest.
